% !TeX root = ../main.tex

\xchapter{结论与展望}{Conclusions and Future Work}

本文围绕 Wi-Fi CSI 室内动作识别中类别集合持续扩展的问题展开研究，分别针对新增类别样本相对充足的类增量学习场景和新增类别样本极少的少样本类增量学习场景开展方法设计与实验验证。全文从任务定义、模型构建和实验分析三个层面出发，重点讨论在旧类数据难以完整保留、新类样本数量受限以及 CSI 时序特征易受环境扰动和类别相似性影响的条件下，如何提升 Wi-Fi 感知模型的持续更新能力。基于前文研究内容，本章对全文主要结论进行归纳，并进一步说明后续可完善的研究方向。

\xsection{结论}{Conclusions}

（1）研究对象与任务界定。本文讨论的并非固定类别集合下的一次性离线训练与识别，而是面向类别持续扩展条件下的 Wi-Fi CSI 室内动作识别：模型需在类别集合随应用逐步扩大的前提下完成更新，并在测试阶段对截至当前的全部已见类别统一判别。围绕上述目标，本文将问题划分为类增量学习与少样本类增量学习两类典型情形，前者对应新增类样本相对充足、可通过记忆库等方式保留旧类信息的设定，后者对应增量阶段每类仅有极少标注、旧类原始训练数据往往难以完整依赖的设定；两类情形体现了从「样本相对充足」到「样本极度稀缺」的递进关系，也为后续方法设计提供了不同的约束界面。

（2）第三章类增量学习方法与实验结论。针对新增类样本相对充足的类增量学习场景，本文提出回放—扩展—巩固框架：通过基于类均值的代表样本选取构建记忆库并完成旧类回放；在冻结 CSI 特征提取骨干的前提下扩展分类器以接入新类别，减弱新类梯度对旧类表征空间的直接破坏；进而采用教师—学生结构与双边知识蒸馏将扩展阶段积累的新旧类知识巩固回单一主干，并结合平衡微调与分类器重训练等环节缓解增量阶段的新旧类样本不平衡。在 XRF 与 WiAR 两个 Wi-Fi CSI 动作识别数据集上的实验表明，该框架在多阶段类别扩展过程中能够缓解灾难性遗忘并保持较好的增量稳定性：在 XRF 上末次增量阶段准确率达到 73.57\%；在 WiAR-16 类增量设置下，平均准确率达到 77.33\%，末次阶段准确率达到 64.58\%。上述结果表明，在可保留有限旧类代表样本的条件下，通过回放、扩展与巩固的衔接，能够在持续接入新类的同时维持整体判别能力。

（3）第四章少样本类增量学习方法与实验结论。针对新增类样本极少的少样本类增量学习场景，本文提出基于时序增强与图注意力的学习方法：利用多尺度时序编码器与时序注意力融合模块强化 CSI 动作序列的时序表征；在增量阶段通过增强图注意力网络建模类别原型之间的关系；并采用双分类器架构融合不同来源的判别信息，以缓解新旧类 logits 偏置与类别级过拟合。在 XRF 的 5-way 1-shot 设定下，本文方法末次增量会话准确率达到 74.24\%，平均准确率达到 78.48\%；在 WiAR 数据集上平均准确率为 75.30\%，末次会话准确率为 61.81\%。这表明在增量监督极为稀疏、难以依赖旧类原始数据的条件下，结合 CSI 时序特性进行表征增强与类别关系建模，仍有助于提升新类适应能力并维持对旧类的判别。

（4）两项研究的递进关系与总体认识。第三章与第四章并非简单并列：第三章侧重于在新增类样本相对充足时，如何在旧类保持与新类接入之间取得平衡并完成跨阶段知识整合；第四章进一步面向新增类样本极少的情形，强调通过时序表征增强与类别关系传播改善原型估计与决策边界。二者共同构成面向 Wi-Fi CSI 动作识别持续更新的一条研究路径。总体认识在于：该类任务除了要抑制灾难性遗忘外，还需结合 CSI 时序结构、类别相似性以及样本规模约束进行针对性设计，单一照搬视觉领域增量设定难以充分刻画无线感知场景下的实际困难。

\xsection{展望}{Future Work}

在本文实验设定下，方法在 XRF、WiAR 等数据集上取得了与对照相比更为稳健的表现，但验证场景仍以固定环境与给定划分为前提，尚未系统覆盖部署环境中的全部变异来源。Wi-Fi CSI 易受房间布局、设备位置与朝向、人员体型与动作习惯以及环境噪声等因素影响，信号分布偏移可能削弱增量阶段的判别边界。在本文工作的基础上，未来可进一步结合域适应、域泛化、元学习或自监督预训练等思路，提升模型在跨环境、跨设备与跨用户条件下的稳健性，并在更严格的跨域协议下开展验证。

第三章框架依赖记忆库存储旧类代表样本以支撑回放与蒸馏，能够在既定存储预算内缓解遗忘；但随着类别数目持续增长，存储占用与隐私合规成本仍可能成为部署瓶颈，也可能限制在敏感场景下的适用性。其原因在于 exemplar 级回放直接保留了原始样本信息。未来可在本文路线之上探索特征级回放、生成式回放、原型约束、参数正则化或无样本蒸馏等机制，在降低对旧类原始样本依赖的同时维持增量阶段的稳定性。

本文主要在离线实验环境中完成模型训练与阶段性评估，而实际 Wi-Fi 感知系统常部署于资源受限的边缘节点，并对延迟与能耗敏感，难以完全复现实验室的批量训练节奏。该局限来自当前研究重心在于方法有效性与可复现对比，而非端到端系统工程。未来可结合模型压缩、参数高效微调、流式数据接入与在线增量更新等手段，在保持增量性能的前提下降低计算与存储开销，使方法更贴近长期在线运行的应用场景。
